\documentclass{article}
\usepackage{amsmath}
\usepackage{amssymb}
\usepackage{array}
\usepackage{booktabs}
\usepackage{textcomp}
\usepackage{graphicx}

\usepackage{amsmath,amsfonts,amsthm,tikz} %packages for math symbols
\newtheorem{theorem}{Theorem}[section]
\newtheorem{lemma}{Lemma}[section]

\newenvironment{solution}
  {\begin{proof}[Solution]}
  {\end{proof}}


\title{diagram for keeping track of terms in part 2}
\author{Ian Chi}
\date{\today}
\begin{document}
\maketitle

\(C\) is a subset that can be removed to get opt. it is of size \(\varepsilon_f n\) where n is the size of the list.\\

\(\delta\) big elements  \(i \in [n]\) and have another \(j\in [n]\) such that \(\frac{1}{2} - \delta\) percent of elements between \(i\) and \(j\) violate monotonicity. 

Claim: opt requires removing exactly half of all 0-big\\


~\\
\quad

\begin{tabular}{ |c|c|c|c| } 
    term & directional & can be in c & must be in c\\
    \hline
    \(\delta\) big      & n & y & n\\
    high/low critical   & y & y & n\\
    left/right big      & y & y & n
\end{tabular}

~\\
\(\delta \) big implies it is either left or right big and vice versa.

\end{document}